\section{Model}
\setlength{\abovedisplayskip}{0pt plus 1pt}%
\setlength{\belowdisplayskip}{0pt plus 1pt}%
\setlength{\abovedisplayshortskip}{0pt plus 1pt}%
\setlength{\belowdisplayshortskip}{0pt plus 1pt}%

To be able to find Skyrmions we first have to define a system in which they might exist. After this some characteristics of these Skyrmions will be discussed. 

\subsection{Hamiltonian}

The system will be an $N$ by $N$ 2D lattice where we have set $N$  to 40. In the sites of this lattice electrons are bound with a spin. The spins of these electrons interact via a normal nearest neighbor ferromagnetic Heisenberg model. A magnetic field is applied therefore a Zeeman Hamiltonian is added. Then to stimulate the underlying crystal structure the Dzyaloshinskii-Mariya (DM) interaction energy is also added.

\begin{equation}
	\mathcal{H} =\mathcal{H}_H + \mathcal{H}_{Z} +\mathcal{H}_{DM}
	\label{eq: H_tot}
\end{equation}


\subsubsection{Heisenberg}

The Heisenberg Hamiltonian describes the alignment between nearest neigbours and is given by: 
\begin{equation}
	\mathcal{H}_H = -J \sum_r \vec{S}_r \cdot \left( \vec{S}_{r + \hat{x}} + \vec{S}_{r + \hat{y}}  \right).
\end{equation}
Here $J$ is the coupling constant that is kept at a constant value of 1, and $\vec{S}_r$ is the spin at site $r$. The sum over $r$ sums over all spins with periodic boundary conditions. 
For positive or negative J it leads to the alignment or anti alignment of nearest neigbours and is therefore ferromagnetic or  antiferromagnetic.


\subsubsection{Zeeman}

The Zeeman Hamiltonian describes the alignment of the spins with an external magnetic field and is described by:

\begin{equation}
	 \mathcal{H}_{Z} = - H_z \sum_r \vec{S}_r 
\end{equation}

Here $H_Z$ is the external magnetic field. Together with the Heisenberg Hamiltonian it describes the Ising model.



\subsubsection{Dzyaloshinskii–Moriya Interaction}
Contrarie to the Ising model the Dzyaloshinskii–Moriya interaction otherwise known as the antisymmetric exchange favors spin canting. 
The phenomena was first observed in weak ferromagnetism of typical antiferromagnetic materials.
The interaction was first described using landau theory by  Igor Dzyaloshinskii in 1958 \cite{Dzyaloshinsky1958ATT} and later using spin orbit coupling by Toru Moriya in 1960 \cite{PhysRev.120.91}.
\begin{equation}
	\mathcal{H}_{DM} = -D \sum_r \vec{S}_r \times \vec{S}_{r + \hat{x}} \cdot \hat{x} + \vec{S}_r \times \vec{S}_{r + \hat{y}} \cdot \hat{y}
\end{equation}
Here $D$ is the DM constant. This interaction makes it possible for nonlinear spin textures to form and therefore makes the formation of skyrmions possible.


\subsubsection{Skyrmion Number}\label{model_skyrmion_number}
Skyrmions are vortices in the spin lattice that can be mathematically described by the skyrmion number otherwise known as the winding number. 

In the continuous case the skyrion number is given by:
\begin{equation}
	 n = \frac{1}{4\pi} \int \mathbf{M} \cdot \left( \frac{\partial \mathbf{M}}{\partial x}  \times \frac{\partial \mathbf{M}}{\partial y}\right) dx dy.
\end{equation}

To calculate this discreetly we use the formula proposed by  Berg and Lüscher (1981)\cite{BERG1981412}:
\begin{equation}
	Q = \frac{1}{4\pi} \sum_{i=1}^N \sum_{j=1}^N \left( q_{ij}^A + q_{ij}^B\right).
\end{equation}
Where $q_{ij}^A $ and $ q_{ij}^B$ are the solid angles of the two triangles in the plaquette (i, j):
\begin{align}
& q_{ij}^A = {} 2\arctan\biggl( \\
& \frac{\mathbf{S}_{i,j}\cdot \left(\mathbf{S}_{i+1,j} \times \mathbf{S}_{i,j+1}\right)}
       {1 + \mathbf{S}_{i,j}\cdot\mathbf{S}_{i+1,j} + \mathbf{S}_{i+1,j}\cdot\mathbf{S}_{i,j+1} + \mathbf{S}_{i,j+1}\cdot\mathbf{S}_{i,j}}
\biggr),  \nonumber \\
& q_{ij}^B = {} 2\arctan\biggl( \\
& \frac{\mathbf{S}_{i+1,j}\cdot \left(\mathbf{S}_{i+1,j+1} \times \mathbf{S}_{i,j+1}\right)}
       {1 + \mathbf{S}_{i+1 ,j}\cdot\mathbf{S}_{i+1,j+1} + \mathbf{S}_{i+1,j+1}\cdot\mathbf{S}_{i,j+1} + \mathbf{S}_{i,j+1}\cdot\mathbf{S}_{i+1,j}}
\biggr) \nonumber
\end{align}
	



